\chapter{Mood Swings}

\section*{Definition}
Rapid or extreme fluctuations in emotional state, shifting between euphoria, irritability, sadness, or anger over short periods without clear precipitants.

\section*{Pathophysiology}
Mood regulation involves complex interplay between neurotransmitters (serotonin, dopamine, norepinephrine), hormones (estrogen, progesterone, cortisol), and neural circuits (prefrontal cortex, amygdala, hippocampus). Disruption from hormonal changes, stress, sleep deprivation, or psychiatric conditions produces emotional lability.

\section*{Etiology}
\begin{itemize}
  \item Premenstrual syndrome (PMS) or PMDD
  \item Perimenopause and menopause
  \item Pregnancy and postpartum period
  \item Bipolar disorder
  \item Depression or anxiety disorders
  \item Chronic stress
  \item Sleep deprivation
  \item Substance use or withdrawal
  \item Thyroid disorders (hyperthyroidism)
  \item Medication side effects (corticosteroids, antidepressants)
\end{itemize}

\section*{Indications for Evaluation}
Seek care if:
\begin{itemize}
  \item Severe or worsening symptoms
  \item Mood swings with thoughts of self-harm, significant functional impairment, or accompanying severe depression or mania
  \item Accompanied by other concerning signs
\end{itemize}

\section*{Related Symptoms}
\begin{itemize}
  \item Irritability
  \item Fatigue
  \item Sleep changes
  \item Anxiety
  \item Appetite changes
\end{itemize}
\textbf{Note:} Always consider serious underlying conditions when evaluating persistent mood swings.

\section*{Self-Care / Home Management}
\begin{itemize}
  \item Rest and avoid activities that may aggravate the condition.
  \item Maintain adequate hydration and a balanced diet.
  \item Over-the-counter remedies may provide symptomatic relief (consult a pharmacist or doctor).
  \item Monitor symptoms and seek professional advice if they persist or worsen.
  \item Avoid self-medicating with prescription medications without medical guidance.
\end{itemize}

\section*{Diagnostic Approach}
\begin{itemize}
  \item A thorough history and physical examination are the first steps in evaluation.
  \item Laboratory tests (blood work, urinalysis) may be ordered based on clinical suspicion.
  \item Imaging studies (X-ray, ultrasound, CT, MRI) may be indicated when structural pathology is suspected.
  \item Specialist referral is arranged when the diagnosis remains uncertain or requires specific expertise.
\end{itemize}

\section*{Treatment / Management Overview}
\begin{itemize}
  \item Treatment targets the underlying cause identified through diagnostic evaluation.
  \item Supportive care and symptom management are provided while awaiting definitive therapy.
  \item Medications, lifestyle modifications, and physical therapies may be prescribed as appropriate.
  \item Follow-up ensures treatment effectiveness and allows timely adjustments to the care plan.
\end{itemize}

\clearpage